\documentclass[11pt,a4paper]{article}
\usepackage[T1]{fontenc}
\usepackage{amsmath,amssymb,geometry}
\geometry{margin=27mm}
\title{EPPA--III: the complete neighbourhood-profile orbit count}
\author{Independent validation supplement; manuscript frozen}
\date{10 October 2026}
\begin{document}
\maketitle

\paragraph{Purpose and source.}
This supplies an exact, self-contained justification of
the replacement proposed for Claim~\texttt{claim:connections}
(Claim~3.9) of \texttt{main-new(7).tex}.
It is independent of any assumption that an EPPA extension
stabilizes an individual complementary clique \(D_j\).
No original manuscript text has been modified.

\paragraph{The orbit formula.}
Let \(G=sK_t\) with components \(C_1,\dots,C_s\).
Fix a vertex \(w\in H\setminus G\) and write
\[
  P(w)=(P_1,\ldots,P_s),\qquad
  P_i=N_H(w)\cap C_i.
\]
For \(0\le j\le t\), put
\(m_j=|\{i:|P_i|=j\}|\), so that
\(\sum_{j=0}^t m_j=s\).
Then
\[
 \boxed{
 \bigl|\operatorname{Aut}(G)\cdot P(w)\bigr|
 =\frac{s!}{\prod_{j=0}^t m_j!}
     \prod_{j=0}^t\binom tj^{m_j}.}
 \tag{1}
\]
Indeed, \(\operatorname{Aut}(G)=S_t\wr S_s\) can
permute the \(s\) clique positions arbitrarily
and can permute vertices inside each clique
independently. Two profiles belong to the
same orbit if and only if their multisets
of coordinate sizes agree. The multinomial
factor chooses which positions have each size
and the second factor independently chooses
the actual subsets at those positions.
This proves (1), not merely an upper bound.

\paragraph{The Claim 3.9 estimate.}
Suppose \(1\le k\le t-1\) and \(m_k=r\ge1\).
Write \(B=\binom tk\ge2\).
By (1), simply choosing the \(r\) positions
whose coordinates have size \(k\), and the
\(k\)-subsets in those positions, gives
\[
 |\operatorname{Aut}(G)\cdot P(w)|
 \ge\binom sr B^r
 \ge sB.
 \tag{2}
\]
For the second inequality, if \(r<s\), use
\(\binom sr\ge s\) and \(B^r\ge B\).
If \(r=s\), use
\(\binom ss=1\) and \(B^{s-1}\ge 2^{s-1}\ge s\).
The last bound follows by induction on \(s\ge1\).

Every automorphism of \(G\) extends to an
automorphism of \(H\). Such an extension
preserves \(G\) and \(H\setminus G\) setwise
and takes \(P(w)\) to its translate.
Even without choosing the extensions coherently,
every profile in (1) is therefore realized
by some vertex of \(H\setminus G\). Consequently,
if \(|H\setminus G|=st\),
\[
          s\binom tk\le st,\qquad
          \binom tk\le t.
\]
For \(2\le k\le t-2\), Pascal's identity
implies \(\binom tk>t\): in the expansion
\(\binom tk=\binom{t-1}{k-1}+
\binom{t-1}{k}\), both terms are at least
\(t-1\). Hence \(k=1\) or \(k=t-1\).

\paragraph{What this does and does not repair.}
The orbit argument itself does not use
regularity of the bipartite graphs
between \(C_i\) and \(D_j\); it proves
an \emph{extreme-degree restriction on
each outside vertex}. To conclude that
each entire bipartite graph is a matching
or co-matching, one still needs the
two-sided regularity asserted in
Claim~\texttt{claim:regularity}, or else
a replacement proof of Lemma~3.7.
The final paragraph of Claim~3.8
is separately flagged as C14 in
\texttt{validation/suggested-corrections.tex}.
It is not silently repaired here.

\paragraph{Independent verification.}
Lean formally proves the general EPPA
upper bound in
\texttt{EPPAIII/EPPA/ProfileOrbitBound.lean}
and the numerical inequalities in
\texttt{EPPAIII/CliqueCovers/Claim39Arithmetic.lean}.
The general orbit \emph{lower bound} (2) is now
proved in Lean in
\texttt{EPPAIII/CliqueCovers/WreathOrbitCount.lean},
using the explicit generators
\texttt{WreathGenerators.lean},
their action on slices
\texttt{WreathSlices.lean},
the independent-subset realizability
\texttt{WreathChoiceTransitivity.lean}, and
the cardinal calculation
\texttt{ProfileChoices.lean}.
The precise \emph{multinomial equality} (1)
remains a paper-level derivation rather
than a Lean theorem; this equality is not
needed for Claim~3.9's numerical conclusion.
An independently generated enumeration
using elementary adjacent permutations,
rather than (1), verifies the count and
bounds for all profiles with \(st\le12\);
see \texttt{checks/section3\_wreath\_orbits.py}.
The checker covers 269 orbits, 19,724
profiles and 257 nontrivial-coordinate
checks.

The final transport from the generated
wreath orbit into the EPPA profile orbit is
prepared in \texttt{Claim39Orbit.lean}
(PR~\#26), still awaiting CI at this
checkpoint. This does not remove the
separate Claim~3.8/C14 regularity dependency
for the entire bipartite blocks.

\end{document}
